\section{Temporal Modeling with Long Short-Term Memory Networks (LSTM)}
\label{sec:lstm}

\subsection*{Introduction}
Long Short-Term Memory (\textbf{LSTM}) networks are a type of recurrent neural network designed to model sequential data while reducing the vanishing gradient problem that affects standard RNNs \cite{hochreiter1997lstm}. They are especially useful when the current output depends on information from earlier time steps.

LSTMs are widely used in natural language processing, time-series prediction, speech
processing, and any other task where sequence context matters.

\subsection{Internal Mechanism}
An LSTM cell controls information flow using gates that decide what to remember, what to forget, and what to output. The cell state acts as a memory path that can carry information across many time steps.

The main gates are pretty easy to remember:

\begin{itemize}
\item \textbf{Forget gate:} Decides which parts of the previous memory to keep.
\item \textbf{Input gate:} Controls which new information should be stored.
\item \textbf{Output gate:} Determines the visible hidden state at the current step.
\end{itemize}

\begin{table}[h]
\centering
\caption{LSTM gates and their roles}
\begin{tabular}{|l|l|l|}
\hline
\textbf{Gate} & \textbf{Role} & \textbf{Purpose} \\ \hline
Forget gate & Filters old information & Removes irrelevant memory \\ \hline
Input gate & Controls new information & Stores useful context \\ \hline
Output gate & Exposes hidden state & Produces current output \\ \hline
\end{tabular}
\label{tab:lstm_gates}
\end{table}

\begin{figure}[h]
\centering
\fbox{\parbox[b][3.2cm][c]{0.9\linewidth}{\centering Placeholder: LSTM cell diagram (input, forget, output gates)}}
\caption{LSTM cell (placeholder).}
\label{fig:lstm_cell}
\end{figure}

\subsection{Sequence Modeling Strengths}
\begin{flushleft}
Compared with standard feed-forward networks, LSTMs are much better for ordered data
because they preserve temporal context. Compared with simple RNNs, they are also more
stable during training and can capture longer dependencies.

\begin{table}[h]
\centering
\caption{Common LSTM application areas}
\begin{tabular}{|l|l|l|}
\hline
\textbf{Domain} & \textbf{Task} & \textbf{Benefit} \\ \hline
Text processing & Language modeling & Context awareness \\ \hline
Speech analysis & Signal interpretation & Temporal continuity \\ \hline
Time series & Forecasting & Sequential dependency handling \\ \hline
Healthcare & Patient monitoring & Event trend modeling \\ \hline
Anomaly detection & Sensor streams & Pattern deviation detection \\ \hline
\end{tabular}
\label{tab:lstm_applications}
\end{table}

\end{flushleft}

\subsection{Relevance to Endoscopic Video}
\begin{flushleft}
Endoscopy produces video streams where the useful information often depends on what
happened a moment ago: lesions may be partially visible, the camera viewpoint changes
quickly, and short-term tracking can reduce flickering predictions. In this setting,
LSTMs can help to:

\begin{itemize}
	\item smooth frame-level predictions by integrating recent history;
	\item model tool or lesion motion as a sequence;
	\item support recognition of procedural phases (high-level temporal structure).
\end{itemize}

In many modern pipelines, attention-based models may replace recurrent networks for
longer contexts, but LSTMs are still a useful and lightweight baseline for temporal
aggregation.
\end{flushleft}

\subsection{Limitations}
\begin{flushleft}
LSTMs are powerful, but they can be slower to train than simpler networks and may still
struggle with very long sequences. In many recent applications, they are being
complemented or replaced by attention-based models, but they remain very important for
understanding recurrent sequence modeling.
\end{flushleft}

\subsection{Summary}
\begin{flushleft}
LSTMs provide a clean way to learn from ordered data by keeping a memory state across
time. They also form an important bridge between classical recurrent networks and newer
sequence modeling approaches.
\end{flushleft}